\section*{第 \S\arabic{exercisegroup} 节习题}
\phantomsection
\addcontentsline{toc}{section}{第 \protect\S\arabic{exercisegroup} 节习题}

\begin{enumerate}
\def\labelenumi{\arabic{enumi})}
\tightlist
\item
  在与 I, p.~20 的命题 2 相同的假设下，证明公式
\end{enumerate}

\[
[ \mathbf {f} ^ {(n)}. \mathbf {g} ] = \sum_ {p = 0} ^ {n} (- 1) ^ {p} \binom {n} {p} D ^ {n - p} [ \mathbf {f}. \mathbf {g} ^ {(p)} ].
\]

\begin{enumerate}
\def\labelenumi{\arabic{enumi})}
\setcounter{enumi}{1}
\tightlist
\item
  采用 I, p.~28 的命题 2 的记号，假设关系 \([\mathbf{a}.\mathbf{y}] = 0\)（对一切 \(\mathbf{y} \in \mathrm{F}\)）蕴涵 E 中的 \(\mathbf{a} = 0\)。在这些条件下，若 \(\mathbf{g}_i\)（\(0 \leqslant i \leqslant n\)）是 \(n + 1\) 个取值于 E 的向量函数，定义在 R 的区间 I 上，并且对每个取值于 F、在 I 上 n 次可导的向量函数 f 都恒有
\end{enumerate}

\[
[ \mathbf {g} _ {0}. \mathbf {f} ] + [ \mathbf {g} _ {1}. \mathbf {f} ^ {\prime} ] + \dots + [ \mathbf {g} _ {n}. \mathbf {f} ^ {(n)} ] = 0
\]

则诸函数 \(g_{i}\) 恒为零。

\begin{enumerate}
\def\labelenumi{\arabic{enumi})}
\setcounter{enumi}{2}
\tightlist
\item
  采用习题 2 的记号并对 \([x.y]\) 作同样的假设，设每个函数 \(g_{k}\) 都在 I 上 n 次可导；对每个在 I 上 n 次可导、取值于 F 的函数 f，令
\end{enumerate}

\[
[ \mathbf {g} _ {0}. \mathbf {f} ] - [ \mathbf {g} _ {1}. \mathbf {f} ] ^ {\prime} + [ \mathbf {g} _ {2}. \mathbf {f} ] ^ {\prime \prime} + \dots + (- 1) ^ {n} [ \mathbf {g} _ {n}. \mathbf {f} ] ^ {(n)} = [ \mathbf {h} _ {0}. \mathbf {f} ] + [ \mathbf {h} _ {1}. \mathbf {f} ^ {\prime} ] + \dots + [ \mathbf {h} _ {n}. \mathbf {f} ^ {(n)} ],
\]

这便无歧义地定义了诸函数 \(\mathbf{h}_i\)（\(0 \leqslant i \leqslant n\)）（习题 2）；证明有

\[
[ \mathbf {h} _ {0}. \mathbf {f} ] - [ \mathbf {h}. \mathbf {f} ] ^ {\prime} + [ \mathbf {h} _ {2}. \mathbf {f} ] ^ {\prime \prime} + \dots + (- 1) ^ {n} [ \mathbf {h} _ {n}. \mathbf {f} ] ^ {(n)} = [ \mathbf {g} _ {0}. \mathbf {f} ] + [ \mathbf {g} _ {1}. \mathbf {f} ^ {\prime} ] + \dots + [ \mathbf {g} _ {n}. \mathbf {f} ^ {(n)} ]
\]

恒成立。

\begin{enumerate}
\def\labelenumi{\arabic{enumi})}
\setcounter{enumi}{3}
\tightlist
\item
  设 \(\mathbf{f}\) 是 \(n\) 次可导的向量函数，定义在区间 \(\mathrm{I} \subset \mathbf{R}\) 上。证明对 \(1 / x \in \mathrm{I}\) 恒有
\end{enumerate}

\[
\frac {1}{x ^ {n + 1}} \mathbf {f} ^ {(n)} \left(\frac {1}{x}\right) = (- 1) ^ {n} \mathrm{D} ^ {n} \left(x ^ {n - 1} \mathbf {f} \left(\frac {1}{x}\right)\right)
\]

（对 \(n\) 作归纳论证）。

\begin{enumerate}
\def\labelenumi{\arabic{enumi})}
\setcounter{enumi}{4}
\tightlist
\item
  设 \(u\) 与 \(v\) 是两个 \(n\) 次可导的实函数，定义在区间 \(I \subset \mathbf{R}\) 上。若令 \(D^n(u / v) = (-1)^n w_n / v^{n+1}\)（在使 \(v(x) \neq 0\) 的每一点处），证明
\end{enumerate}

\[
w _ {n} = \left| \begin{array}{c c c c c c} u & v & 0 & 0 & \dots & 0 \\ u ^ {\prime} & v ^ {\prime} & v & 0 & \dots & 0 \\ u ^ {\prime \prime} & v ^ {\prime \prime} & 2 v ^ {\prime} & v & \dots & 0 \\ \dots & \dots & \dots & \dots & \dots & \dots \\ u ^ {(n - 1)} & v ^ {(n - 1)} & \binom {n - 1} {1} v ^ {(n - 2)} & \binom {n - 1} {2} v ^ {(n - 3)} & \dots & v \\ u ^ {(n)} & v ^ {(n)} & \binom {n} {1} v ^ {(n - 1)} & \binom {n} {2} v ^ {(n - 2)} & \dots & \binom {n} {n - 1} v ^ {\prime} \end{array} \right|
\]

（令 \(w = u / v\) 并微分 \(n\) 次关系 \(u = wv\)）。

\begin{enumerate}
\def\labelenumi{\arabic{enumi})}
\setcounter{enumi}{5}
\tightlist
\item
  设 \(\mathbf{f}\) 是定义在开区间 \(\mathrm{I} \subset \mathbf{R}\) 上、取值于赋范空间 \(\mathrm{E}\) 的向量函数。
\end{enumerate}

令 \(\Delta \mathbf{f}(x; h_1) = \mathbf{f}(x + h_1) - \mathbf{f}(x)\)，然后归纳地定义

\[
\Delta^ {p} \mathbf {f} (x; h _ {1}, h _ {2}, \dots , h _ {p - 1}, h _ {p}) = \Delta^ {p - 1} \mathbf {f} (x + h _ {p}; h _ {1}, \dots , h _ {p - 1}) - \Delta^ {p - 1} \mathbf {f} (x; h _ {1}, \dots , h _ {p - 1});
\]

这些函数对每个 \(x \in I\) 都有定义（当诸 \(h_{i}\) 充分小时）。

\begin{enumerate}
\def\labelenumi{\alph{enumi})}
\tightlist
\item
  若函数 \(\mathbf{f}\) 是 \(n\) 次可导的（在点 \(x\) 处；从而是 \(n - 1\) 次可导的，在 \(x\) 的一个邻域上），则有
\end{enumerate}

\[
\lim_{\substack{(h_{1},\ldots ,h_{n})\to (0,\ldots ,0)\\ h_{1}h_{2}\ldots h_{n}\neq 0}}\frac{\varDelta^{n}\mathbf{f}(x;h_{1},\ldots,h_{n})}{h_{1}h_{2}\ldots h_{n}} = \mathbf{f}^{(n)}(x)
\]

（对 n 作归纳论证，并使用中值定理）。

\begin{enumerate}
\def\labelenumi{\alph{enumi})}
\setcounter{enumi}{1}
\tightlist
\item
  若 \(\mathbf{f}\) 在区间 I 上 \(n\) 次可导，则有
\end{enumerate}

\[
\begin{array}{l} \left\| \Delta^ {n} \mathbf {f} (x; h _ {1}, \ldots , h _ {n}) - \mathbf {f} ^ {(n)} (x _ {0}) h _ {1} h _ {2} \ldots h _ {n} \right\| \\ \leqslant | h _ {1} h _ {2} \ldots h _ {n} | \sup \left\| \mathbf {f} ^ {(n)} (x + t _ {1} h _ {1} + \dots + t _ {n} h _ {n}) - \mathbf {f} ^ {(n)} (x _ {0}) \right\| \end{array}
\]

其中上确界是对 \((t_i)\)（满足 \(0 \leqslant t_i \leqslant 1\)，对 \(1 \leqslant i \leqslant n\)）所成的集取的（同样方法）。

\begin{enumerate}
\def\labelenumi{\alph{enumi})}
\setcounter{enumi}{2}
\tightlist
\item
  若 \(f\) 是在 I 上 \(n\) 次可导的实函数，则有
\end{enumerate}

\[
\Delta^ {n} f (x; h _ {1}, h _ {2}, \dots , h _ {n}) = h _ {1} h _ {2} \dots h _ {n} f ^ {(n)} (x + \theta_ {1} h _ {1} + \dots + \theta_ {n} h _ {n})
\]

其中诸数 \(\theta_{i}\) 属于 {[}0, 1{]}（同样方法，利用 I, p.~22 的推论）。

\begin{enumerate}
\def\labelenumi{\arabic{enumi})}
\setcounter{enumi}{6}
\tightlist
\item
  设 \(f\) 是 \(n\) 次可导的实值有限函数（在点 \(x_0\) 处），\(\mathbf{g}\) 是 \(n\) 次可导的向量函数（在点 \(y_0 = f(x_0)\) 处）。设
\end{enumerate}

\[
\begin{array}{l} f (x _ {0} + h) = a _ {0} + a _ {1} h + \dots + a _ {n} h ^ {n} + r _ {n} (h) \\ \mathbf {g} (y _ {0} + k) = \mathbf {b} _ {0} + \mathbf {b} _ {1} k + \dots + \mathbf {b} _ {n} k ^ {n} + \mathbf {s} _ {n} (k) \end{array}
\]

分别是 f 在点 \(x_{0}\) 处与 g 在点 \(y_{0}\) 处的 n 阶泰勒展开。证明 \(n+1\) 项之和（指复合函数 \(g \circ f\) 在点 \(x_{0}\) 处的 n 阶泰勒展开中的）等于下列多项式中次数 \(\leqslant n\) 的项之和

\[
\mathbf {b} _ {0} + \mathbf {b} _ {1} (a _ {1} h + \dots + a _ {n} h ^ {n}) + \mathbf {b} _ {2} (a _ {1} h + \dots + a _ {n} h ^ {n}) ^ {2} + \dots + \mathbf {b} _ {n} (a _ {1} h + \dots + a _ {n} h ^ {n}) ^ {n}.
\]

由此推出下面两个公式：

\begin{enumerate}
\def\labelenumi{\alph{enumi})}
\tightlist
\item
\end{enumerate}

\[
\mathrm{D} ^ {n} (\mathbf {g} (f (x))) = \sum \frac {n !}{m _ {1} ! m _ {2} ! \dots m _ {q} !} \mathbf {g} ^ {(p)} (f (x)) \left(\frac {f ^ {\prime} (x)}{1 !}\right) ^ {m _ {1}} \dots \left(\frac {f ^ {(q)} (x)}{q !}\right) ^ {m _ {q}}
\]

其中求和取遍一切满足如下条件的正整数组 \((m_{i})_{1\leqslant i\leqslant q}\)：

\[
m _ {1} + 2 m _ {2} + \dots + q m _ {q} = n
\]

其中 p 表示和 \(m_{1} + m_{2} + \cdots + m_{q}\)。

\begin{enumerate}
\def\labelenumi{\alph{enumi})}
\setcounter{enumi}{1}
\tightlist
\item
\end{enumerate}

\[
\mathrm{D} ^ {n} (\mathbf {g} (f (x))) = \sum_ {p = 1} ^ {n} \frac {1}{p !}   \mathbf {g} ^ {(p)} (f (x)) \left(\sum_ {q = 1} ^ {p} \binom {p} {q} (- f (x)) ^ {p - q}   \mathrm{D} ^ {n} ((f (x)) ^ {q}\right).
\]

\begin{enumerate}
\def\labelenumi{\arabic{enumi})}
\setcounter{enumi}{7}
\tightlist
\item
  设 \(f\) 是一个 \(n\) 次可导、定义在区间 \(\mathrm{I}\) 上的实函数，设 \(x_{1}, x_{2}, \ldots, x_{p}\) 是 \(\mathrm{I}\) 的互异点，而 \(n_{i}\)（\(1 \leqslant i \leqslant p\)）是 \(p\) 个整数 \(> 0\)，使得
\end{enumerate}

\[
n _ {1} + n _ {2} + \dots + n _ {p} = n.
\]

假设在点 \(x_{i}\) 处，函数 f 连同它的前 \(n_{i}-1\) 阶导数对 \(1 \leqslant i \leqslant p\) 都为零：证明存在一点 \(\xi\)，它内在于包含诸 \(x_{i}\) 的最小区间，并且使得 \(f^{(n-1)}(\xi) = 0\)。

\begin{enumerate}
\def\labelenumi{\arabic{enumi})}
\setcounter{enumi}{8}
\tightlist
\item
  采用与习题 8 相同的记号，设 \(f\) 在 I 上 \(n\) 次可导但别处任意。设 \(g\) 是 \(n - 1\) 次（实系数）多项式，使得在点 \(x_{i}\)（\(1 \leqslant i \leqslant p\)）处，\(g\) 及其前 \(n_{i} - 1\) 阶导数分别等于 \(f\) 及其前 \(n_{i} - 1\) 阶导数。证明我们有
\end{enumerate}

\[
f (x) = g (x) + \frac {(x - x _ {1}) ^ {n _ {1}} (x - x _ {2}) ^ {n _ {2}} \dots (x - x _ {p}) ^ {n _ {p}}}{n !} f ^ {(n)} (\xi)
\]

其中 \(\xi\) 内在于包含诸点 \(x_{i}\)（\(1 \leqslant i \leqslant p\)）与 \(x\) 的最小区间。（对 \(t\) 的函数

\[
f (t) - g (t) - a \frac {(t - x _ {1}) ^ {n _ {1}} (t - x _ {2}) ^ {n _ {2}} \dots (t - x _ {p}) ^ {n _ {p}}}{n !}
\]

应用习题 8，其中 \(a\) 是适当选取的常数。）

\begin{enumerate}
\def\labelenumi{\arabic{enumi})}
\setcounter{enumi}{9}
\tightlist
\item
  设 \(g\) 是定义在 0 的一个邻域上的奇实函数，并且在该邻域上 5 次可导。证明
\end{enumerate}

\[
g (x) = \frac {x}{3} \left(g ^ {\prime} (x) + 2 g ^ {\prime} (0)\right) - \frac {x ^ {5}}{1 8 0} g ^ {(5)} (\xi) \quad (\xi = \theta x, \quad 0 <   \theta <   1)
\]

（与习题 9 相同的方法）。

由此推出：若 \(f\) 是定义在 \([a, b]\) 上且在此区间上 5 次可导的实函数，则

\[
f (b) - f (a) = \frac {b - a}{6} \left[ f ^ {\prime} (a) + f ^ {\prime} (b) + 4 f ^ {\prime} \left(\frac {a + b}{2}\right) \right] - \frac {(b - a) ^ {5}}{2 8 8 0} f ^ {(5)} (\xi)
\]

其中 \(a < \xi < b\)（“辛普森公式”）。

\begin{enumerate}
\def\labelenumi{\arabic{enumi})}
\setcounter{enumi}{10}
\tightlist
\item
  设 \(f_{1}, f_{2}, \ldots, f_{n}\) 与 \(g_{1}, g_{2}, \ldots, g_{n}\) 是 \(2n\) 个在区间 I 上 \(n - 1\) 次可导的实函数。设 \((x_{i})_{1 \leqslant i \leqslant n}\) 是 I 中点的严格递增序列。证明两个行列式
\end{enumerate}

\[
\left| \begin{array}{c c c c} f _ {1} (x _ {1}) & f _ {1} (x _ {2}) & \dots & f _ {1} (x _ {n}) \\ f _ {2} (x _ {1}) & f _ {2} (x _ {2}) & \dots & f _ {2} (x _ {n}) \\ \dots & \dots & \dots & \dots \\ f _ {n} (x _ {1}) & f _ {n} (x _ {2}) & \dots & f _ {n} (x _ {n}) \end{array} \right| \vdots \left| \begin{array}{c c c c} g _ {1} (x _ {1}) & g _ {1} (x _ {2}) & \dots & g _ {1} (x _ {n}) \\ g _ {2} (x _ {1}) & g _ {2} (x _ {2}) & \dots & g _ {2} (x _ {n}) \\ \dots & \dots & \dots & \dots \\ g _ {n} (x _ {1}) & g _ {n} (x _ {2}) & \dots & g _ {n} (x _ {n}) \end{array} \right|
\]

之比等于两个行列式

\[
\left| \begin{array}{c c c c} f _ {1} (\xi_ {1}) & f _ {1} ^ {\prime} (\xi_ {2}) & \dots & f _ {1} ^ {(n - 1)} (\xi_ {n}) \\ f _ {2} (\xi_ {1}) & f _ {2} ^ {\prime} (\xi_ {2}) & \dots & f _ {2} ^ {(n - 1)} (\xi_ {n}) \\ \dots & \dots & \dots & \dots \\ f _ {n} (\xi_ {1}) & f _ {n} ^ {\prime} (\xi_ {2}) & \dots & f _ {n} ^ {(n - 1)} (\xi_ {n}) \end{array} \right|: \left| \begin{array}{c c c c} g _ {1} (\xi_ {1}) & g _ {1} ^ {\prime} (\xi_ {2}) & \dots & g _ {1} ^ {(n - 1)} (\xi_ {n}) \\ g _ {2} (\xi_ {1}) & g _ {2} ^ {\prime} (\xi_ {2}) & \dots & g _ {2} ^ {(n - 1)} (\xi_ {n}) \\ \dots & \dots & \dots & \dots \\ g _ {n} (\xi_ {1}) & g _ {n} ^ {\prime} (\xi_ {2}) & \dots & g _ {n} ^ {(n - 1)} (\xi_ {n}) \end{array} \right|
\]

之比，其中

\[
\xi_ {1} = x _ {1}, \quad \xi_ {1} <   \xi_ {2} <   x _ {2}, \quad \xi_ {2} <   \xi_ {3} <   x _ {3}, \quad \dots , \quad \xi_ {n - 1} <   \xi_ {n} <   x _ {n}
\]

（应用 I, p.~38 的习题 9）。

\(g_{1}(x)=1\)、\(g_{2}(x)=x,\ldots,g_{n}(x)=x^{n-1}\) 的特殊情形。

¶ 12) \(a)\) 设 \(\mathbf{f}\) 是定义在有限区间 \(\mathrm{I} = [-a, + a]\) 上且连续的向量函数，取值于赋范空间 \(\mathrm{E}\)（在 \(\mathbf{R}\) 上）中，并且在 \(\mathrm{I}\) 上两次可导。若令 \(\mathrm{M}_0 = \sup_{x\in \mathrm{I}}\| \mathbf{f}(x)\|\)、\(\mathrm{M}_2 = \sup_{x\in \mathrm{I}}\| \mathbf{f}''(x)\|\)，证明对一切 \(x\in \mathrm{I}\) 有

\[
\left\| \mathbf {f} ^ {\prime} (x) \right\| \leqslant \frac {\mathrm{M} _ {0}}{a} + \frac {x ^ {2} + a ^ {2}}{2 a} \mathrm{M} _ {2}
\]

（把诸差 \(\mathbf{f}(a) - \mathbf{f}(x)\)、\(\mathbf{f}(-a) - \mathbf{f}(x)\) 表出）。

\begin{enumerate}
\def\labelenumi{\alph{enumi})}
\setcounter{enumi}{1}
\tightlist
\item
  由 a) 推出：若 f 是区间 I（有界或否）上的两次可导函数，并且 \(\mathbf{M}_{0} = \sup_{x \in I} \| \mathbf{f}(x) \|\) 与 \(\mathbf{M}_{2} = \sup_{x \in I} \| \mathbf{f}''(x) \|\) 有限，则 \(\mathbf{M}_{1} = \sup_{x \in I} \| \mathbf{f}'(x) \|\) 也有限，并且有：
\end{enumerate}

\[
\mathrm{M} _ {1} \leqslant 2 \sqrt {\mathrm{M} _ {0} \mathrm{M} _ {2}} \quad \text { if   I   has   length } \geqslant 2 \sqrt {\frac {M _ {0}}{\mathrm{M} _ {2}}}
\]

\[
\mathrm{M} _ {1} \leqslant \sqrt {2} \sqrt {\mathrm{M} _ {0} \mathrm{M} _ {2}} \quad \text { if } \mathrm{I} = \mathbf {R}.
\]

证明在这两个不等式中，数 2 与 \(\sqrt{2}\) 分别不能换成更小的数（先考虑仅设 f 具有二阶右导数的情形，并证明此时取 f 为“分段”等于二次多项式的实函数，前述不等式中的两项可以变成相等的）。

\begin{enumerate}
\def\labelenumi{\alph{enumi})}
\setcounter{enumi}{2}
\tightlist
\item
  由 \(b)\) 推出：若 \(\mathbf{f}\) 是 \(p\) 次可导的（在 \(\mathbf{R}\) 上），并且 \(\mathbf{M}_p = \sup_{x \in \mathbf{R}} \| \mathbf{f}^{(p)}(x) \|\)
\end{enumerate}

与 \(M_{0}=\sup_{x\in\mathbf{R}}\|\mathbf{f}(x)\|\) 有限，则数 \(M_{k}=\sup_{x\in\mathbf{R}}\left\|\mathbf{f}^{(k)}(x)\right\|\) 中的每一个都有限（因为

\(1 \leqslant k \leqslant p - 1)\) 以及

\[
\mathbf {M} _ {k} \leqslant 2 ^ {k (p - k) / 2} \mathbf {M} _ {0} ^ {1 - k / p} \mathbf {M} _ {p} ^ {k / p}.
\]

¶ 13) a) 设 \(f\) 是 \(\mathbf{R}\) 上的两次可导实函数，使得 \((f(x))^2 \leqslant a\) 与 \((f'(x))^2 + (f''(x))^2 \leqslant b\) 在 \(\mathbf{R}\) 上成立；证明

\[
(f (x)) ^ {2} + \left(f ^ {\prime} (x)\right) ^ {2} \leqslant \max (a, b)
\]

在 R 上成立（用反证法论证，注意若函数 \(f^{2}+f^{\prime2}\) 取值 \(c > \max(a, b)\)（在某点 \(x_{0}\) 处），则存在两点 \(x_{1}, x_{2}\)，使得 \(x_{1} < x_{0} < x_{2}\)，并且在 \(x_{1}\) 与 \(x_{2}\) 处函数 \(f^{\prime}\) 取的值足够小，使 \(f^{2} + f^{\prime2}\) 取值 \textless{} c；然后考虑 \([x_{1}, x_{2}]\) 中使 \(f^{2} + f^{\prime2}\) 在该区间上达到其上确界的一点）。

\begin{enumerate}
\def\labelenumi{\alph{enumi})}
\setcounter{enumi}{1}
\tightlist
\item
  设 \(f\) 是 \(n\) 次可导的实函数，定义在 \(\mathbf{R}\) 上，使得 \((f(x))^2 \leqslant a\) 与 \((f^{(n - 1)}(x))^2 + (f^{(n)}(x))^2 \leqslant b\) 在 \(\mathbf{R}\) 上成立；证明这时
\end{enumerate}

\[
(f ^ {(k - 1} (x)) ^ {2} + (f ^ {(k)} (x)) ^ {2} \leqslant \max (a, b)
\]

在 \(\mathbf{R}\) 上对 \(1 \leqslant k \leqslant n\) 成立。（对 \(n\) 作归纳论证；注意，依习题 12，上确界 \(c\)（\((f'(x))^2\) 在 \(\mathbf{R}\) 上的）是有限的；通过归结为矛盾来证明必有 \(c \leqslant \max(a, b)\)：假设 \(c > \max(a, b)\)，选取常数 \(\lambda\) 与 \(\mu\)，使得对函数 \(g = \lambda f + \mu\) 有 \(|g(x)| \leqslant 1\)、\(|g'(x)| \leqslant 1\)，然而 \((g(x))^2 + (g'(x))^2 \leqslant 1\) 不可能对一切 \(x\) 都成立。

¶ 14) 设 \(\mathbf{f}\) 是在包含 0 的区间 I 上 \(n - 1\) 次可导的函数，并设 \(\mathbf{f}_n\) 是对 \(x \neq 0\) 由关系

\[
\mathbf {f} (x) = \mathbf {f} (0) + \mathbf {f} ^ {\prime} (0) \frac {x}{1 !} + \mathbf {f} ^ {\prime \prime} (0) \frac {x ^ {2}}{2 !} + \dots + \mathbf {f} ^ {(n - 1)} (0) \frac {x ^ {n - 1}}{(n - 1) !} + \mathbf {f} _ {n} (x) x ^ {n}.
\]

\begin{enumerate}
\def\labelenumi{\alph{enumi})}
\item
  证明若 \(\mathbf{f}\) 在点 0 处有 \((n + p)^{th}\) 阶导数，则 \(\mathbf{f}_n\) 在点 0 处有 \(p^{th}\) 阶导数，并且在 0 的一个邻域中异于 0 的所有点处有 \((n + p - 1)^{th}\) 阶导数；此外，\(\mathbf{f}_n^{(k)}(0) = \frac{k!}{(n + k)!}\mathbf{f}^{(n + k)}(0)\) 对 \(0 \leqslant k \leqslant p\) 成立，并且 \(\mathbf{f}_n^{(p + k)}(x)x^k\) 随 \(x\) 趋于 0（对 \(1 \leqslant k \leqslant n - 1\)）（把 \(\mathbf{f}_n\) 的导数借助 \(\mathbf{f}\) 的各阶导数的泰勒展开表出，并使用 I, p.~18 的命题 6）。
\item
  反之，设 \(f_{n}\) 是在 I 中 0 的一个邻域上具有 \((n + p - 1)^{th}\) 阶导数的向量函数，并且 \(\mathbf{f}_{n}^{(p+k)}(x)x^{k}\) 对 \(0 \leqslant k \leqslant n - 1\) 有极限。证明函数 \(\mathbf{f}_{n}(x)x^{n}\) 在 I 上有 \((n + p - 1)^{th}\) 阶导数；若进一步 \(f_{n}\) 在点 0 处具有 \(p^{th}\) 阶导数，则 \(\mathbf{f}_{n}(x)x^{n}\) 在点 0 处具有 \((n + p)^{th}\) 阶导数。
\item
  假设 I 关于 0 对称且 f 是偶函数（在 I 上 \(\mathbf{f}(-x) = \mathbf{f}(x)\)）。借助 a) 证明：若 f 在 I 上 2n 次可导，则存在定义在 I 上且 n 次可导的函数 g，使得在 I 上 \(\mathbf{f}(x) = \mathbf{g}(x^{2})\)。
\end{enumerate}

¶ 15) 设 I 是 \(\mathbf{R}\) 中的开区间，\(\mathbf{f}\) 是定义在 I 上且连续的向量函数；假设存在 \(n\) 个定义在 I 上的向量函数 \(\mathbf{g}_i\)（\(1 \leqslant i \leqslant n\)），使得 \(x\) 的函数

\[
\frac {1}{h ^ {n}} \left(\mathbf {f} (x + h) - \mathbf {f} (x) - \sum_ {p = 1} ^ {n} \frac {h ^ {p}}{p !} \mathbf {g} _ {p} (x)\right)
\]

当 h 趋于 0 时在含于 I 的每个紧区间上一致趋于 0。

\begin{enumerate}
\def\labelenumi{\alph{enumi})}
\item
  我们令 \(\mathbf{f}_p(x,h) = \Delta^p\mathbf{f}(x;h,h,\ldots ,h)\)（I, p.~40, 习题 6）。证明：对 \(1\leqslant p\leqslant n\)，\((1 / h^{p})\mathbf{f}_{p}(x,h)\) 在 I 的每个紧子区间上一致趋于 \(\mathbf{g}_p(x)\)（当 \(h\) 趋于 0 时），并且诸 \(\mathbf{g}_p\) 在 I 上连续（依次对 \(p = n\)、\(p = n - 1\) 等证明这一点）
\item
  由此推出 \(\mathbf{f}\) 具有连续的 \(n^{th}\) 阶导数，并且 \(\mathbf{f}^{(p)} = \mathbf{g}_p\) 对 \(1 \leqslant p \leqslant n\) 成立（把关系 \(\mathbf{f}_{p+1}(x, h) = \mathbf{f}(x + h, h) - \mathbf{f}_p(x, h)\) 考虑在内）。
\end{enumerate}

¶ 16) 设 \(f\) 是 \(n\) 次可导的实函数，定义在 \(\mathrm{I} = ] - 1, + 1[\) 上，并且在此区间上 \(|f(x)|\leqslant 1\)。

\begin{enumerate}
\def\labelenumi{\alph{enumi})}
\tightlist
\item
  证明若 \(m_k(\lambda)\) 表示 \(|f^{(k)}(x)|\) 在含于 I 的长度为 \(\lambda\) 的区间上的最小值，则有
\end{enumerate}

\[
m _ {k} (\lambda) \leqslant \frac {2 ^ {k (k + 1) / 2} k ^ {k}}{\lambda^ {k}}
\]

\[
(1 \leqslant k \leqslant n).
\]

（注意若把长度为 \(\lambda\) 的区间分解为长度为 \(\alpha, \beta, \gamma\) 的三个区间，则有

\[
m _ {k} (\lambda) \leqslant \frac {1}{\beta} (m _ {k - 1} (\alpha) + m _ {k - 1} (\gamma)).
\]

\begin{enumerate}
\def\labelenumi{\alph{enumi})}
\setcounter{enumi}{1}
\tightlist
\item
  由 a) 推出：存在仅依赖于整数 n 的数 \(\mu_{n}\)，使得若 \(|f'(0)| \geqslant \mu_{n}\)，则导数 \(f^{(n)}(x)\) 在 I 的至少 n - 1 个互异的点上为零（用对 k 的归纳法证明 \(f^{(k)}\) 在 I 上至少零 k - 1 次）。
\end{enumerate}

\begin{enumerate}
\def\labelenumi{\arabic{enumi})}
\setcounter{enumi}{16}
\tightlist
\item
  \begin{enumerate}
  \def\labelenumii{\alph{enumii})}
  \tightlist
  \item
    设 f 是开区间 \(I \subset R\) 上具有各阶导数的向量函数。假设在 I 上有 \(\left\|\mathbf{f}^{(n)}(x)\right\| \leqslant a n!r^n\)，其中 a 与 r 是两个 \textgreater0 且与 x 和 n 无关的数；证明在每个点 \(x_0\) 处，通项为 \((1/n!) \mathbf{f}^{(n)}(x_0)(x - x_0)^n\) 的“泰勒级数”收敛，并且其和为 \(\mathbf{f}(x)\)（在 \(x_0\) 的某个邻域上）。
  \end{enumerate}
\end{enumerate}

\begin{enumerate}
\def\labelenumi{\alph{enumi})}
\setcounter{enumi}{1}
\item
  反之，若 \(\mathbf{f}\) 在点 \(x_0\) 处的泰勒级数在 \(x_0\) 的某个邻域上收敛，则存在两个数 \(a\) 与 \(r\)（依赖于 \(x_0\)），使得 \(\left\| \mathbf{f}^{(n)}(x_0)\right\| \leqslant a.n!r^n\) 对每个整数 \(n > 0\) 成立。
\item
  由 a) 与习题 16 b) 推出：若在开区间 \(\mathrm{I} \subset \mathbf{R}\) 上实函数 \(f\) 无穷次可导，并且存在整数 \(p\)（与 \(n\) 无关），使得对一切 \(n\)，函数 \(f^{(n)}\) 至多有 \(p\) 个在 \(\mathrm{I}\) 中不为零的互异点，则 \(f\) 在每个点 \(x_0 \in \mathrm{I}\) 的一个邻域上的泰勒级数收敛，并且其和为 \(f(x)\)（在 \(x_0\) 的一个邻域的每个点处）。
\end{enumerate}

\begin{enumerate}
\def\labelenumi{\arabic{enumi})}
\setcounter{enumi}{17}
\tightlist
\item
  设 \((a_{n})_{n\geqslant 0}\) 是任意一个复数序列。对每个 \(n\geqslant 0\)，令 \(s_n^{(0)} = a_n\)，并且对 \(k\geqslant 0\) 归纳地定义
\end{enumerate}

\[
s _ {n} ^ {(k + 1)} = s _ {0} ^ {(k)} + s _ {1} ^ {(k)} + \dots + s _ {n - 1} ^ {(k)}.
\]

\begin{enumerate}
\def\labelenumi{\alph{enumi})}
\tightlist
\item
  证明“序列的泰勒公式”：对每个整数
\end{enumerate}

\[
\left| s _ {n + h} ^ {(k)} - s _ {n} ^ {(k)} - h s _ {n} ^ {(k - 1)} - \binom {h} {2} s _ {n} ^ {(k - 2)} - \dots - \binom {h} {k - 1} s _ {n} ^ {(1)} \right| \leqslant \binom {h} {k} \sup _ {0 \leqslant j \leqslant h - 1} | a _ {n + j} |
\]

（对 k 施行归纳）。

\begin{enumerate}
\def\labelenumi{\alph{enumi})}
\setcounter{enumi}{1}
\tightlist
\item
  假设存在数 C，使得对一切 n 有 \(|na_{n}| \leqslant C\)，并且序列 \((s^{(2)}/n)\)（由算术平均 \((s_{0} + \cdots + s_{n-1})/n\)、即部分和 \(s_{n} = a_{0} + \cdots + a_{n-1}\) 的算术平均组成）趋于极限 \(\sigma\)。证明通项为 \(a_{n}\) 的级数收敛且其和为 \(\sigma\)（“哈代–李特尔伍德陶伯定理”）。（写
\end{enumerate}

\[
s _ {n} = \frac {1}{h} \left(s _ {n + h} ^ {(2)} - s ^ {(2)}\right) + \frac {h - 1}{2} r _ {n}
\]

其中 \(|r_n|\) 借助不等式 \(|na_n| \leqslant C\) 有上界，而 \(h\) 作为 \(n\) 的函数适当选取。）

\setcounter{exercisegroup}{3}
\refstepcounter{exercisegroup}
